% Part: many-valued-logic
% Chapter: infinite-valued-logics
% Section: lukasiewicz

\documentclass[../../../include/open-logic-section]{subfiles}

\begin{document}

\olfileid{mvl}{inf}{luk}

\olsection{\L ukasiewicz తర్కం}

\begin{editorial}
  ఇది అనంత-విలువల \L ukasiewicz తర్కంపై
  సంక్షిప్త ప్రాథమిక విభాగం మాత్రమే.
\end{editorial}

\begin{defn}\ollabel{def:lukasiewicz}
అనంత-విలువల \L ukasiewicz
తర్కం~$\LogLuk[\infty]$ను కింది మాత్రికతో
నిర్వచిస్తాం:
\begin{enumerate}
  \item $\lfalse$, $\lnot$, $\land$, $\lor$, $\lif$
  సంయోజకాలు గల ప్రామాణిక ప్రతిజ్ఞావాక్య
  భాష $\Lang L_0$.
  \item సత్యమూల్యాల సమితి $V_\infty$.
  \item $1$ మాత్రమే నిర్దేశిత విలువ; అంటే
  $V^+ = \{1\}$.
  \item సత్యమూల్య ప్రమేయాలను కింది
  ప్రమేయాలు ఇస్తాయి:
  \begin{align*}
    \tf{\lfalse} & = 0\\
    \tf{\lnot}[\LogLuk](x) & = 1 - x\\
    \tf{\land}[\LogLuk](x,y) & = \min(x,y)\\
    \tf{\lor}[\LogLuk](x,y) & = \max(x,y)\\
    \tf{\lif}[\LogLuk](x,y) & = \min(1,1-(x-y)) = \begin{cases}
      1 & \text{ఒకవేళ } x \le y\\
      1-(x-y) & \text{లేకపోతే.}
    \end{cases}
    \end{align*}
\end{enumerate}
$V = V_m$ తీసుకోవడం మినహా $m$-విలువల
\L ukasiewicz తర్కాన్ని ఇదే విధంగా
నిర్వచిస్తాం.
\sourcecorrection{OLTEMVLINF-003}{ఆంగ్ల
మూలం ప్రామాణిక భాషను పేర్కొన్నా,
దానిలోని అసత్య స్థిరాంకాన్ని జాబితాలో
గానీ ప్రమేయాల్లో గానీ ఇవ్వలేదు.
మూడు-విలువల నిర్వచనంతో తరువాతి
సమానత్వ వాదం సరిపోవడానికి అసత్య
స్థిరాంకాన్ని చేర్చి దానికి సున్నా
విలువను సంపాదకీయంగా ఇచ్చాం. నాలుగు
ముద్రిత ప్రమేయాలు మాత్రమే ఆ విలువను
నిర్బంధించవు.}
\end{defn}

\begin{prop}
  \olref[thr][luk]{def:lukasiewicz}లో
  నిర్వచించిన $\LogLuk[3]$ తర్కమే
  \olref{def:lukasiewicz}లో
  నిర్వచించిన $\LogLuk[3]$ తర్కం.
\end{prop}

\begin{proof}
  \olref[thr][luk]{def:lukasiewicz}లో
  ఇచ్చిన సంయోజకాల సత్య పట్టికలను,
  \olref{def:lukasiewicz}లోని
  సమీకరణాలు నిర్ణయించే కింది సత్య
  పట్టికలతో పోల్చి దీన్ని చూడవచ్చు:
  \begin{center}
    \begin{tabular}{c|c}
      $\tf{\lnot}$ & \\
      \hline
      $1$ & $0$ \\
      $1/2$ & $1/2$ \\
      $0$ & $1$
    \end{tabular}
    \quad
    \begin{tabular}{c|ccc}
      $\tf{\land}[\LogLuk[3]]$ & $1$ & $1/2$ & $0$ \\
      \hline
      $1$ & $1$ & $1/2$ & $0$ \\
      $1/2$ & $1/2$ & $1/2$ & $0$\\
      $0$ & $0$ & $0$ & $0$
    \end{tabular}
    \\[2ex]
    \begin{tabular}{c|ccc}
      $\tf{\lor}[\LogLuk[3]]$ & $1$ & $1/2$ & $0$ \\
      \hline
      $1$ & $1$ & $1$ & $1$ \\
      $1/2$ & $1$ & $1/2$ & $1/2$ \\
      $0$ & $1$ & $1/2$ & $0$
    \end{tabular}
    \quad
    \begin{tabular}{c|ccc}
      $\tf{\lif}[\LogLuk[3]]$ & $1$ & $1/2$ & $0$ \\
      \hline
      $1$ & $1$ & $1/2$ & $0$ \\
      $1/2$ & $1$ & $1$ & $1/2$  \\
      $0$ & $1$ & $1$ & $1$
    \end{tabular}
  \end{center}
\end{proof}

\begin{prop}\ollabel{prop:luk-infty-m}
  $\Gamma \Entails[\LogLuk[\infty]] !B$
  అయితే ప్రతి~$m \ge 2$కు
  $\Gamma \Entails[\LogLuk[m]] !B$.
\end{prop}

\begin{proof}
  సాధన.
\end{proof}

\begin{prob}
  \olref[mvl][inf][luk]{prop:luk-infty-m}ను
  నిరూపించండి.
\end{prob}

నిజానికి తిరుగు దిశ కూడా నిలుస్తుంది.

అనంత-విలువల \L ukasiewicz తర్కం
అత్యంత ప్రసిద్ధ ఫజీ తర్కం.
ఫజీ తర్క సాహిత్యంలో సోపాధికాన్ని
తరచూ $\lnot !A \lor !B$గా
నిర్వచిస్తారు. అలా చేస్తే
అనంత-విలువల బలమైన క్లీని తర్కం
లభిస్తుంది.

\begin{prob}
  $(p \lif q) \lor (q \lif p)$
  ~$\LogLuk[\infty]$లో సర్వసత్యమని
  చూపండి.
\end{prob}

\end{document}
